\ssscp{Putative PCEMP *kera(nŋ) `hawksbill turtle'}{13-01-05-03}
Elsewhere \citep[46]{Blust2009},
PCEMP *keRa(nŋ) ‘hawksbill turtle: \it{Chelonia imbricata}’\footnote{
		Modern sources use \it{Eretmochelys imbricata}, not \it{Chelonia
		imbricata}, for ‘hawksbill turtle’.}
is given as providing additional
support for PCEMP *e.
As \citet[45]{Blust2009} interprets the data,
“The known evidence suggests that PAN/PMP *peñu may have referred
to both the green turtle
and the hawksbill, and that these referents may not have been
lexically distinguished until AN speakers reached eastern Indonesia.”
At that point he sees a lexical development that adds PCEMP *keRa(nŋ)
‘hawksbill turtle: \it{Chelonia imbricata}’ that coexisted with
\tsc{PAn} *pəñu ‘sea turtle’,
with the PMP gloss of ‘the green turtle, \it{Chelonia mydas’}.
PCEMP *keRa(nŋ) “has a surprisingly light
and scattered distribution” \citeyear[46]{Blust2009}.\footnote{
		In a related
		story of the lexicon,
		\tsc{PAn} *qaCipa is glossed ‘soft-shelled turtle’,
		and PMP *qantipa is glossed ‘type of freshwater turtle’.
		In spite of there being no reflexes in any language glossed ‘hawksbill
		turtle’ for these etyma
		and glosses such as ‘type of river turtle’,
		\citet{BlustTrussel2020} speculate in a note that it is “likely
		that PAn *qaCipa and PMP *qantipa referred to the ocean-dwelling hawksbill turtle,
		which is prized for its widely-traded shell,
		than to the economically unimportant land tortoise.
		If so, *qaCipa was replaced by *keRa/keRa(nŋ) in PCEMP.”}

We see the data telling a different story,
beginning far to the west of the Blust line.
For Indonesian (Malay), 
\citet{StevensSchmidgall-Tellings2010} list the form \it{kəraŋ}
‘shell; generic name for cockles,
oysters, clams, mussels’.\footnote{
		Although \citet{StevensSchmidgall-Tellings2010} label \it{kəraŋ} ‘obsolete’,
		nevertheless in September 2020,
		several 20-year-old university students in eastern Indonesia
		(Kupang) supplied \it{kəraŋ} as their first response for generic
		‘seashell’, so there appear to be regional variations of usage.}
\citet{StevensSchmidgall-Tellings2010} also note that the associated
term \it{kéraŋ} [keraŋ] ‘tortoise shell’ is used in Indonesian
in eastern Indonesia.\footnote{
		Orthographic ‹é› in Indonesian
		dictionaries indicates /e/ rather than /ə/.
		The two phonemes are underdifferentiated in the Indonesian orthography.
		Words with /e/ are low frequency.}
Most varieties of contact Malay
in eastern Indonesia have only five vowels,
and regularly shift penultimate schwa *ə from the western varieties
of Malay to \it{e} or \it{a},
depending on the region.
We note that in our experience \it{kulit kéraŋ} (with \it{e}
and initial stress; [ˈkeraŋ]) or \it{kulit karáŋ} (with \it{a}
and ultimate stress; [kaˈraŋ])\footnote{
		Such \it{e}
		and \it{a} variants are common for Malay words with penultimate
		schwa in the west.
		As a result, \it{karáŋ} ‘shell’ with ultimate stress is a minimal
		pair with \it{káraŋ} ‘coral’ with initial stress.}
are widely
used in Malay varieties around eastern Indonesia to refer to
‘any kind of shell from the sea’,
with the exception of Ambon Malay,
which uses \it{kulit bia} ‘any kind of shell from the sea’.
Kupang Malay \it{kéoŋ} [ˈkeoŋ] ‘seashell’ is often heard alongside
\it{karáŋ} (cf. Malay \it{kéoŋ} ‘snail’).

\citet{EcholsWolff1989} similarly list Indonesian \it{kəraŋ}
‘cockle shells, mollusks and their shell’; \it{bər-kəraŋ} ‘cultivate shellfish’;
\it{kəraŋ{\tl}kəraŋ-an} ‘various kinds of mollusks’; \it{pə-kəraŋ-an} ‘heap of shells’.

\citet{Nothofer1992} pointed out other possible cognates in MP
languages west of the Blust line such as Padoe \it{kéha} ‘shell’,
and Kada Toraja \it{kéra} ‘shell’.
\citet[46]{Blust2009} dismisses this saying
“[\ldots] if these forms are related, they almost certainly are loan-
words, since (1) the Toraja languages are spoken in interior Sulawesi, and their speakers
would have no knowledge of the hawksbill turtle apart, perhaps, from its valued shell,
which would have been acquired in trade, and (2) as already noted, the vowel *e (not *ə)
has not been reconstructed for \tsc{PAn} or PMP.”
 


However, Blust’s response ignores several issues.
When Grimes collected data in the Padoe region in South Sulawesi,
that language was spoken both in the mountains
and along the coast,
so there is no reason to assume their word for “shell” should
be a loan (see map in \citealp[20]{GrimesGrimes1987}).
Furthermore, Grimes has encountered numerous seashells (including
several varieties of clams,
conches, top shells, and a chambered nautilus) at lake Rana in
the mountainous centre of Buru Island -- four-days’ difficult
walk from the coast in any direction.
When queried what they were doing with the seashells,
the response was “We brought them up to slake lime for chewing betelnut.
We didn’t have time to do that when we were down at the coast.”
(See \frf{13-02}.) People in the interior of Buru know what season
of the year and what phase of the moon is best to go down to the coast
and get \it{peno telu-n} ‘sea turtle eggs’.
They spend a large percentage of their time as hunters and gatherers.
They often travel for days
and spend extended time away from their villages to follow the
food sources, such as during the east monsoon extended hunting season.
Land snails are often described with terms that relate them to seashells.
So Buru \it{feku} ‘land snail with shell’ (< PMP *bəkul ‘snail’)
is also described as a \it{kima mua} ‘jungle mollusc’ (\it{kima}
‘any seashell or the creature that lives in them’ < PMP *kima
‘giant clam’; see \citealp{GrimesGrimes2020}).
So it should not be assumed that Austronesian-speaking people
in the mountains have loanwords for seashells or do not know
about sea turtles apart from trade.

\begin{figure}[h]\centering
	\includegraphics[width=0.5\columnwidth]{13-02.jpg}
	\fcp{Burning seashells in the mountains of Buru to slake lime for chewing betelnut}{13-02}
\end{figure}

\citet[vol. 2; 827]{Wolff2010} relates the terms associated with various
types of ‘turtle’ in the CEMP region (including Oceania) to a
PMP form *gaɣaŋ ‘crab’ with the note “In l[an]g[uage]s S[outh]
of the Phil[ippines], this root developed *e in the 1st syl[la]b[lle]
and the meaning shifted from ‘crab’ to ‘shellfish’.” Wolff’s
data show that the gloss “shell” is not only found in languages
west of the Blust line such as Moken \it{ʔiek klāŋ} ‘the ark
shell’, and Malay \it{kəraŋ} ‘cockleshell’,
but the gloss also persisted out into Oceanic languages,
such as Tolai \it{kelaŋ} ‘shell’.
In Blust’s transcription Wolff’s *gaɣaŋ would be written as *kaRaŋ,
although no such form (or similar) is found at the level of PMP
in \citet{BlustTrussel2020}.

Returning to the topic of turtles,
in our experience with both speakers of Austronesian languages
and Australian Aboriginal people in Arnhem Land,
the main focus is on the eggs and the meat.
Once the meat is taken,
practical use of the shells is limited by the inconvenience of
fending off ants, biodegradability of biological material,
and that too much drying in the sun causes the scales to warp
and flake off (see \frf{13-03}).
So the shells tend to be cured carefully only if the intent is
to sell them (to Chinese merchants).
More durable shells, such as the giant clam shell (\it{Tridacna
spp.}), are favoured for practical uses such as for drinking troughs
for animals or for evaporating sea salt.
The species of sea turtle is rarely distinguished,
most often just referring to “its eggs”,
“its meat”, “its skin/shell”.
When the different species are in focus the divisions have many variations.
So Buru distinguishes between \it{gofot} ‘land turtle’
and \it{peno} ‘sea turtle’.
Asilulu adds modifiers, as in \it{wenu} ‘sea turtle’,
\it{wenu iʔan} ‘smooth-backed sea turtle’,
\it{wenu sisir} ‘scaley-backed sea turtle (commercial turtle shell; cf.
Malay \it{sisir} ‘comb’)’,
\it{tipopo} ‘smooth-backed land turtle’ \citep{Collins2003a}.

\begin{figure}[h]\centering
	\includegraphics[width=0.75\columnwidth]{13-03.jpg}
	\fcp{Sea turtle taken for its meat, abandoning its shell}{13-03}
\end{figure}

We note that in the data available for non-Oceanic languages
in the CEMP region,
only one language explicitly associates the term with the ‘hawksbill turtle’.
That is the Leti \it{kerna} ‘middle-sized marine turtle with
prized shell: \it{Chelonia imbricata}’ attributed to Jonker.\footnote{
		The expanded gloss of the Leti term is attributed by \citet[243, 252]{Blust1993} to \citet{Jonker1932}.
		\citet[123]{Jonker1932} only gives ‘het schildpad,
		karet’ = ‘turtle/tortoise, tortoiseshell’.
		Nevertheless, \citet[417]{vanEngelenhoven2004} does give a more
		specific gloss with \it{kɛrna} ‘kind of turtle (hawk’s bill,
		\it{Chelonia imbricata 1.})’}
Other languages vary considerably,
including glosses such as ‘turtle,
turtle shell, land turtle,
small turtle, large turtle,
turtle with prized shell’.
So with such widely varied glosses,
it is by no means clear why the gloss for this reconstruction
should be narrowed specifically to ‘hawksbill turtle’ east of
the Blust line (and west of the EMP line).
The form and gloss of the data in Wallacea
and further west should not be determined by the Oceanic data.

There are forms which are formally
and semantically similar to putative *keRa(nŋ),
such as Proto-Philippine *karáy ‘plastron,
turtle shell’ and Proto-Philippine *kara ‘tortoise shell (?)’.
Dempwolff reconstructed *karaq ‘shell’,
which included Tagalog \it{kála} ‘tortoise,
tortoise shell’, Toba Batak \it{hara} ‘turtle’,
and Malagasy \it{hára} ‘mother-of-pearl’.

So how are we to interpret all this diverse data? Because of
differences in the penultimate vowel,
it seems that there was a pair of related forms in western languages
at a stage prior to reaching eastern Indonesia (PMP?): *kəRa(ŋ)
‘shell, turtle shell’ and *kara ‘shell, turtle shell’.
The semantic association of ‘shell’ with ‘turtle shell’ is natural
(as witnessed independently by English, among others).

In eastern Indonesia PMP *ə > \it{e} is extremely common (see
\trf{13-03}), so the forms in most languages can derive directly from *kəRa(ŋ),
as seen in \trf{13-04}.
Potential problem correspondences are shaded.
With the exception of the forms with problems in the sound correspondences,
all can be derived from *kəRaŋ, with schwa.

\begin{table}\tcp{Selected reflexes of *kəRaŋ ‘shell, turtle shell’}{13-04}
\begin{adjustbox}{width=\textwidth}
	\begin{threeparttable}\st{2pt}
	\begin{tabular}{lllllll}\lsptoprule
Language	&	Form	&	Gloss		&	*ə /\und{σ}σ{\#}\su{†}	&	*-R-	&	*-ŋ{\#}	\\	\midrule
Manggarai	&	hera	&	sea turtle			&	ə, e	&	-r-	&	-ŋ	{\cdr}\\	
Ngadha	&	kera	&	large turtle			&	ə, e	&	-r/{\0}-	&	-{\0}	\\	
Palu{\Q}e	&	kea	&	small turtle on reef		(irr., probably loan)	&	ə, (e?)	{\cdr}&	-r/ʤ-	{\cdr}&	-{\0}	\\	
Lio	&	kera	&	tortoise, land turtle			&	ə	{\cdr}&	-r/{\0}-	&	-{\0}	\\	
Sika	&	era	&	turtle			&	ə, e	&	-r-	&	-ŋ	{\cdr}\\	
Havu	&	ea	&	turtle, turtle shell			&	ə, e, o	&	-{\0}-	&	-{\0}	\\	
Termanu, Korbafo, Bokai,	&	kea	&	turtle/tortoise			&	e	&	-{\0}-	&	-{\0}	\\	
Bilbaa, Landu, Ba{\Q}a, Tii			&		&		&		&		&		\\	
(\tsc{Nuclear Rote})			&		&		&		&		&		\\	
Rikou (Rote)	&	ʔea	&	turtle/tortoise		(sporadic *a > \it{e} /{\gap}{\#})	&	e	&	-{\0}-	&	-{\0}	\\	
Dela, Oenale, Dengga	&	ʔee	&	turtle		(sporadic *a > \it{e} /{\gap}{\#})	&	e	&	-{\0}-	&	-{\0}	\\	
(\tsc{West \tsc{Rote-Meto}})	&		&				&		&		&		\\	
Roi{\Q}is Amarasi (\tsc{Meto})	&	kee	&	turtle		(sporadic *a > \it{e} /{\gap}{\#})	&	e	&	-{\0}-	&	-{\0}	\\	
Kotos Amarasi (\tsc{Meto})	&	kea/kee	&	(sea) turtle, (land) tortoise			&	e	&	-{\0}-	&	-{\0}	\\	
Leti	&	kɛrna	&	hawk’s bill turtle 		(C met.)	&	e	&	-r-	&	-n	\\	
Fordata	&	eran	&	turtle with prized shell			&	e, a	&	-r-	&	-n	\\	
Kei	&	keran	&	turtle with prized shell			&	e	&	-r/{\0}-	&	-n	\\	
Kowiai	&	eran	&	turtle			&	e, o, u	&	-r-	&	-n	\\	
		\lspbottomrule
	\end{tabular}
		\begin{tablenotes}\tnsize
			\item[†] {*ə /\und{σ}σ{\#} is schwa in penultimate syllables.}
			%\item[‡]
			%\item[Δ]
			%\item[‖]
			%\item[¶]
		\end{tablenotes}
	\end{threeparttable}
\end{adjustbox}
\end{table}

Moving south and eastward from the Philippines,
the narrowing of the glosses in isolated languages from a fairly
broad \it{shell} ←→ \it{turtle shell > turtle with prized shell
> hawksbill turtle} begins to be seen in just a few languages
near the EMP boundary,
but the broader gloss of ‘shell’ also appears to persist out
into Oceania, such as in Tolai.
The gloss of ‘hawksbill turtle’ is inappropriately narrow to
account for the range of glosses in eastern Indonesia
and the persistence of the gloss for ‘shell’
and ‘small turtle’ into Oceania.

A potential problem with our proposed *kəRaŋ occurs when looking
at Oceanic languages whose reflexes have \it{e},
since PMP *ə > \tsc{POc} *o \citep{Ross1998}.
The solution to that is to move the form with *e further downstream
to \tsc{POc} *keRa(nŋ) ‘shell,
turtle shell, hawksbill turtle: \it{Eretmochelys imbricata}’.
Since Kowiai has been shown to be a \tsc{Seram-Tanimbar-Bomberai}
language, rather than SHWNG (\crf{ch:STB}),
it appears that \tsc{POc} would be the appropriate level for
irregular *ə > \it{e} since we are aware of no cognates from SHWNG languages.

To recapitulate, the evidence for PCEMP *e was originally based
on developments in two words,
to which a third was added.
Blust has since conceded that for *i-sei ‘who?’ convergence is
“difficult to rule out”,
and regarding *ma-qetəm ‘black’,
that he should “not include it as evidence for CEMP.” *keRa(nŋ)
‘hawksbill turtle’ is shown to not be exclusive to the CEMP region,
to be associated with formally close words meaning ‘shell,
turtle shell’ in the Philippines,
western Indonesia and beyond from *kəRa(ŋ)
and formally close words,
with glosses meaning ‘shell’
and ‘small turtle’ persisting into Oceania,
and that the form with *e should be shifted to \tsc{POc}.

Mid-vowels in Wallacean subgroups are from many different sources and pathways.
The occasional appearance of /o/ or /e/ in some of the same words
in different languages in different subgroups (such as Buru
and Selaru) should not be surprising.
But we have shown that this is unpatterned,
and that they got there through different pathways (see footnotes
4 and 5 in this chapter).
As yet, we have found no innovations in the mid-vowels that support
linking some Wallacean subgroups together at a higher level,
and also no support from the mid-vowels for CMP or CEMP.

Note the steps in the arguments for mid-vowels in PCEMP.
Of the four etyma originally proposed \citep{Blust1993} in support
of PCEMP *o and *e, PCEMP *todan ‘sit’, *ma-qetəm ‘black’
and *i-sei ‘who?’ were later withdrawn \citep{Blust2009} as concessions to various criticisms.
PCEMP *oliq ‘return’ has also been shown to be from parallel
development and unnecessary -- a possibility that Blust foresaw.
But perhaps as a later oversight,
in \citet[740]{Blust2013}, PCEMP *todan ‘sit’, *oliq ‘return’,*ma-qetəm
‘black’ and *i-sei ‘who?’ reappear as evidence supporting CEMP mid-vowels,
and why CEMP “is as well-established as any large AN subgroup
apart from Oceanic” \citep[740]{Blust2013}.

Also note methodologically, that lexical evidence from two to
three items for *o
and two to three items for *e,
all of which can be shown to have become similar through parallel
development and are also found in Sulawesi (where mid vowels have also been
reconstructed) outside the CEMP region,
became the basis for proposing a typological innovation in the
phonology from a four-vowel system in PMP to a six-vowel system
for PCEMP, with the addition of the two mid-vowels.
And while most claims about the development of mid-vowels based
on those lexical items were withdrawn in the face of criticisms
from various sources, nevertheless the six-vowel system for PCEMP has been maintained
as if it were well established.
The presence of the mid-vowels in a parent language may only
be appropriate at the level of \tsc{POc}.
